\chapter{Heart Transplant}

\section*{What Is This Surgery?}
Heart transplantation is a complex surgical procedure that involves replacing a patient's severely diseased or failing heart with a healthy heart from a deceased donor. This definitive treatment is reserved for individuals suffering from end-stage heart failure, a condition where the heart can no longer pump enough blood to meet the body's needs, despite maximal medical and device therapies. The procedure primarily targets the thoracic cavity, specifically the mediastinum, where the recipient's native heart is surgically removed and the donor heart is meticulously anastomosed to the great vessels and atrial remnants. It represents a life-saving intervention, significantly improving both the survival rates and quality of life for carefully selected patients who would otherwise face imminent mortality from their cardiac condition. The surgery is not only a technical challenge but also requires a multidisciplinary approach to manage the complex immunological and physiological changes that occur post-transplant.

\section*{Alternative Names}
\begin{itemize}
  \item Cardiac Transplantation
  \item Orthotopic Heart Transplantation (OHT)
  \item Heart Replacement Surgery
  \item Allogeneic Heart Transplantation
  \item Donor Heart Transplant
\end{itemize}

\section*{Relevant Surgical Anatomy}
The heart is centrally located within the thoracic cavity, nestled in the mediastinum, anterior to the vertebral column and posterior to the sternum. It is enclosed within the pericardium, a double-layered sac that provides protection and lubrication. Key anatomical structures relevant to heart transplantation include the four chambers of the heart (right and left atria, right and left ventricles) and the great vessels that enter and exit them. These include the superior vena cava (SVC) and inferior vena cava (IVC) draining into the right atrium, the pulmonary artery originating from the right ventricle, the four pulmonary veins draining into the left atrium, and the aorta originating from the left ventricle.

During orthotopic heart transplantation, the recipient's diseased heart is excised, leaving behind cuffs of the right and left atria, and transected segments of the aorta and pulmonary artery. The donor heart is then implanted, requiring precise anastomoses of the donor's great vessels and atrial remnants to the recipient's corresponding structures. The coronary arteries, which supply blood to the heart muscle, are entirely contained within the donor heart and are thus denervated post-transplant, leading to an absence of typical angina symptoms in the event of coronary allograft vasculopathy. The lymphatic drainage of the heart is primarily to mediastinal lymph nodes, which play a role in the immune response to the transplanted organ. Understanding these intricate relationships is crucial for successful surgical execution and managing postoperative complications.

\section*{Surgical Principle \& Rationale}
The fundamental principle of heart transplantation is the restoration of normal cardiac function and systemic perfusion by replacing a non-functional heart with a healthy one. This is achieved through an orthotopic technique, where the recipient's diseased heart is excised, and the donor heart is implanted in its anatomical position. The surgical rationale centers on correcting the underlying pathophysiology of end-stage heart failure, such as severely impaired ventricular contractility, intractable arrhythmias, or uncorrectable structural defects, which lead to inadequate cardiac output and systemic congestion.

The technical goals involve establishing secure and patent anastomoses between the donor heart and the recipient's great vessels. This typically includes the aorta, the pulmonary artery, and the left and right atria (or superior and inferior venae cavae in the bicaval technique). Once these connections are complete and the donor heart is reperfused, it is expected to resume normal contractile function, thereby restoring adequate blood flow to all organs and resolving the symptoms of heart failure. A critical, ongoing aspect of the principle is the management of the recipient's immune response to the foreign donor organ, which necessitates lifelong immunosuppressive therapy to prevent acute and chronic rejection, ensuring the long-term viability of the transplanted heart.

\section*{History \& Surgical Evolution}
The concept of heart transplantation has roots in early 20th-century experimental surgery. The first successful animal heart transplant was performed by Vladimir Demikhov in the Soviet Union in 1946, demonstrating the feasibility of the procedure in dogs. Building on this pioneering work, the first human heart transplant was famously performed by Dr. Christiaan Barnard in Cape Town, South Africa, on December 3, 1967, on Louis Washkansky. Although Washkansky survived only 18 days, succumbing to pneumonia, this event marked a monumental milestone in medical history.

Following Barnard's breakthrough, the Stanford University team, led by Dr. Norman Shumway and Dr. Richard Lower, refined the orthotopic transplantation technique and established rigorous patient selection criteria, significantly improving early outcomes. However, widespread adoption was limited by the formidable challenge of immune rejection. The true turning point came in the early 1980s with the introduction of cyclosporine, a potent immunosuppressive drug. Cyclosporine dramatically reduced the incidence and severity of acute rejection episodes, transforming heart transplantation from an experimental procedure with dismal long-term survival into an established, life-saving therapy. Subsequent advancements in immunosuppression, donor organ preservation, and mechanical circulatory support devices have further improved outcomes and expanded the indications for heart transplantation in the modern era.

\section*{Clinical Indications}
Heart transplantation is indicated for patients with end-stage heart disease who meet specific criteria and have exhausted other treatment options.

\begin{itemize}
  \item End-stage heart failure (New York Heart Association Class III or IV) with severe symptoms refractory to optimal medical therapy and device management.
  \item Severe ischemic cardiomyopathy with extensive myocardial damage not amenable to surgical revascularization or percutaneous coronary intervention.
  \item Dilated cardiomyopathy with progressive ventricular dysfunction and intractable symptoms despite guideline-directed medical therapy.
  \item Restrictive or hypertrophic cardiomyopathy leading to severe diastolic or systolic dysfunction and debilitating symptoms.
  \item Life-threatening ventricular arrhythmias (e.g., recurrent ventricular tachycardia or fibrillation) refractory to antiarrhythmic medications, catheter ablation, or implantable cardioverter-defibrillator (ICD) therapy.
  \item Inoperable congenital heart disease with uncorrectable structural defects resulting in progressive heart failure.
  \item Severe valvular heart disease with irreversible ventricular dysfunction despite surgical valve repair or replacement.
  \item Primary cardiac tumors that are unresectable and cause significant cardiac dysfunction or metastatic risk.
  \item Refractory cardiogenic shock requiring continuous inotropic support or mechanical circulatory support as a bridge to transplant.
\end{itemize}

\section*{Clinical Positioning}
Heart transplantation serves as the definitive, primary treatment for end-stage heart disease, offering the best long-term survival and quality of life for eligible patients. However, it is typically considered a secondary intervention in the overall management pathway, initiated only after maximal medical therapy, lifestyle modifications, and device-based treatments (such as implantable cardioverter-defibrillators, cardiac resynchronization therapy, or ventricular assist devices) have failed to adequately control symptoms or improve prognosis. In this context, it acts as a salvage procedure for patients whose disease has progressed beyond the reach of conventional therapies.

For some patients, particularly those with rapidly deteriorating conditions or acute cardiogenic shock, mechanical circulatory support devices like ventricular assist devices (VADs) may be implanted as a "bridge to transplant," making the transplant procedure secondary to the VAD implantation. In rare cases of fulminant myocarditis or severe congenital heart disease presenting acutely, heart transplantation may be considered a more primary, urgent intervention if conventional therapies are immediately futile.

\section*{Surgical Technique \& Procedure}
Heart transplantation is a highly complex surgical procedure performed under general anesthesia, typically lasting four to six hours. The patient is positioned supine, and a median sternotomy incision is made to access the heart. Cardiopulmonary bypass (CPB) is then initiated, diverting the patient's blood from the heart and lungs to an external heart-lung machine, which oxygenates the blood and circulates it throughout the body. This allows the surgical team to operate on a still, bloodless heart.

The key operative steps generally follow this sequence:

\begin{itemize}
  \item **Anesthesia and Access:** General anesthesia is induced, and the patient is intubated. Central venous and arterial lines are placed for monitoring. A median sternotomy is performed to expose the heart and great vessels.
  
  \item **Cardiopulmonary Bypass Initiation:** Heparin is administered, and cannulas are placed in the aorta and venae cavae to establish cardiopulmonary bypass, allowing the heart to be arrested and excised.
  
  \item **Recipient Cardiectomy:** The recipient's diseased heart is carefully excised. In the standard orthotopic technique, the great vessels (aorta and pulmonary artery) are transected, and the atria are divided, leaving cuffs of the recipient's right and left atria.
  
  \item **Donor Heart Preparation:** The donor heart, which has been carefully preserved in cold cardioplegic solution and transported, is prepared by trimming the great vessels and atria to match the recipient's anatomy.
  
  \item **Anastomoses:** The donor heart is implanted by meticulously connecting the vessels in a specific order:
  
    \begin{enumerate}
      \item **Left Atrial Anastomosis:** The donor's left atrium is connected to the recipient's left atrial cuff, ensuring proper alignment of the pulmonary veins.
      \item **Right Atrial Anastomosis (Standard Biatrial):** The donor's right atrium is connected to the recipient's right atrial cuff, incorporating the superior and inferior venae cavae.
      \item **Bicaval Anastomosis (Alternative):** In the bicaval technique, the donor's superior vena cava is directly anastomosed to the recipient's SVC, and the donor's inferior vena cava to the recipient's IVC, creating a more anatomically correct right atrium.
      \item **Pulmonary Artery Anastomosis:** The donor pulmonary artery is connected to the recipient's pulmonary artery.
      \item **Aortic Anastomosis:** The donor aorta is connected to the recipient's aorta.
    \end{enumerate}
  
  \item **De-airing and Weaning from CPB:** After all anastomoses are complete, the heart is carefully de-aired to prevent air emboli, and the patient is gradually weaned from cardiopulmonary bypass as the donor heart begins to function.
  
  \item **Hemostasis and Closure:** Once the donor heart is functioning independently, meticulous hemostasis is achieved. Temporary pacing wires are often placed on the epicardium, and chest tubes are inserted to drain any fluid or blood from the mediastinum. The sternum is then closed with stainless steel wires, and the overlying tissues and skin are sutured in layers.
\end{itemize}

\section*{Preoperative Workup \& Preparation}
Extensive preparation is required for heart transplantation, beginning with a comprehensive pre-transplant evaluation to determine eligibility and optimize the patient's condition. This multi-disciplinary assessment includes a thorough medical history, physical examination, and a battery of diagnostic tests.

\begin{itemize}
  \item **Cardiac Evaluation:** Echocardiography, cardiac catheterization (right and left heart) to assess hemodynamics and pulmonary vascular resistance (PVR), stress testing, and sometimes cardiac MRI.
  
  \item **Pulmonary Evaluation:** Pulmonary function tests (PFTs) and chest imaging (X-ray, CT scan).
  
  \item **Renal and Hepatic Function:** Comprehensive blood tests to assess kidney and liver health, including creatinine, BUN, liver enzymes, and bilirubin.
  
  \item **Infectious Disease Screening:** Extensive testing for viral infections (HIV, Hepatitis B \& C, CMV, EBV, VZV), bacterial, and fungal infections, and tuberculosis.
  
  \item **Immunological Assessment:** HLA typing and panel reactive antibody (PRA) levels to assess the risk of rejection and guide donor matching.
  
  \item **Psychosocial Evaluation:** Assessment of mental health, coping mechanisms, social support, and commitment to lifelong adherence to complex medical regimens.
  
  \item **Nutritional Assessment:** Optimization of nutritional status, often involving dietary counseling.
  
  \item **Cancer Screening:** Age-appropriate cancer screening (e.g., colonoscopy, mammography, prostate-specific antigen).
\end{itemize}

Patients are required to cease smoking and illicit drug use, maintain a healthy weight, and receive necessary vaccinations. Once listed, patients await a suitable donor heart, often supported by a ventricular assist device (VAD) if their condition deteriorates. Prior to surgery, patients must be NPO (nothing by mouth) for a specified period, and certain medications, particularly anticoagulants, are adjusted or held. Preoperative antibiotics are administered to minimize infection risk, and detailed informed consent is obtained, outlining the risks, benefits, and alternatives to transplantation.

\section*{Contraindications \& Precautions}
Heart transplantation is a major procedure with significant risks, and careful patient selection is paramount.

Absolute Contraindications:

\begin{itemize}
  \item Active systemic infection or uncontrolled chronic infection (e.g., osteomyelitis, active tuberculosis).
  \item Irreversible pulmonary hypertension (Pulmonary Vascular Resistance > 5 Wood units or transpulmonary gradient > 15 mmHg), which would lead to acute right heart failure in the donor heart.
  \item Active malignancy or a history of malignancy with a high risk of recurrence (typically within 5 years, depending on cancer type and stage).
  \item Severe, irreversible end-organ damage (e.g., renal failure requiring dialysis, severe hepatic cirrhosis, irreversible neurological damage) not amenable to combined organ transplantation.
  \item Active substance abuse (alcohol, illicit drugs) or documented non-compliance with medical regimens.
  \item Severe peripheral or cerebrovascular disease that would preclude safe surgery or recovery.
  \item Lack of adequate psychosocial support system or demonstrated inability to adhere to complex post-transplant care.
\end{itemize}

Relative Contraindications and Precautions:

\begin{itemize}
  \item Advanced age (typically >70 years, though this is institution-dependent and based on physiological age and comorbidity burden).
  \item Significant comorbidities such as poorly controlled diabetes with end-organ damage, severe obesity (BMI >35-40 kg/m\textasciicircum2), or severe osteoporosis.
  \item High panel reactive antibodies (PRA) indicating a high risk of hyperacute or accelerated acute rejection, requiring desensitization protocols.
  \item Previous extensive thoracic surgery, which can complicate re-entry, increase surgical risk, and prolong operative time.
  \item Active peptic ulcer disease or diverticulitis, which can be exacerbated by immunosuppression.
  \item Severe frailty or cachexia, indicating poor physiological reserve.
  \item Psychiatric illness that would impair adherence to post-transplant care, requiring careful evaluation and management.
\end{itemize}

\section*{Complications \& Risks}
Heart transplantation carries significant risks, both immediate and long-term, due to the complexity of the surgery and the necessity of lifelong immunosuppression.

General Surgical Risks:

\begin{itemize}
  \item **Bleeding:** Significant blood loss requiring transfusions, and potentially re-exploration for hemostasis, especially in patients with prior cardiac surgery or on anticoagulants.
  \item **Infection:** Surgical site infection, pneumonia, urinary tract infection, or sepsis, particularly in an immunosuppressed state.
  \item **Anesthesia Complications:** Adverse reactions to anesthetic agents, stroke, myocardial infarction, respiratory failure, or aspiration pneumonia.
\end{itemize}

Procedure-Specific Risks:

\begin{itemize}
  \item **Primary Graft Dysfunction (PGD):** Failure of the donor heart to function adequately immediately after transplantation, often requiring mechanical circulatory support (e.g., ECMO) or re-transplantation.
  \item **Acute Rejection:** The recipient's immune system attacks the donor heart, requiring intensified immunosuppression. Can be cellular or antibody-mediated, diagnosed by endomyocardial biopsy.
  \item **Chronic Rejection (Coronary Allograft Vasculopathy - CAV):** Progressive, diffuse narrowing of the coronary arteries of the donor heart, leading to ischemia and heart failure, often without typical angina due to denervation.
  \item **Immunosuppression-Related Complications:**
    \begin{itemize}
      \item **Opportunistic Infections:** Increased susceptibility to viral (CMV, EBV, HSV, BK virus), fungal (Aspergillus, Candida), and bacterial infections.
      \item **Renal Dysfunction:** Nephrotoxicity from calcineurin inhibitors (e.g., cyclosporine, tacrolimus) leading to chronic kidney disease or failure, often requiring dialysis or kidney transplant.
      \item **Hypertension:** Common side effect requiring aggressive management.
      \item **Diabetes Mellitus:** New-onset diabetes after transplant (NODAT) or worsening of pre-existing diabetes.
      \item **Osteoporosis:** Increased risk of bone loss and fractures due to corticosteroids and other immunosuppressants.
      \item **Malignancy:** Increased risk of certain cancers, particularly skin cancers (squamous cell carcinoma), post-transplant lymphoproliferative disease (PTLD), and Kaposi's sarcoma.
    \end{itemize}
  \item **Arrhythmias:** Bradycardia (due to denervation), atrial fibrillation, or ventricular arrhythmias, often requiring pacing or antiarrhythmic medications.
  \item **Stroke or Neurological Complications:** Due to emboli, hypoperfusion during surgery, or immunosuppressant side effects (e.g., tremors, seizures).
  \item **Gastrointestinal Complications:** Peptic ulcers, esophagitis, or colitis, often exacerbated by medications.
  \item **Death:** Despite advancements, heart transplantation remains a high-risk procedure with a significant mortality rate, especially in the first year.
\end{itemize}

\section*{Postoperative Recovery Timeline}
Immediately following heart transplantation, the patient is transferred to the intensive care unit (ICU) for close monitoring.

\begin{itemize}
  \item **First 24-48 Hours (ICU):** The patient will typically be intubated and mechanically ventilated, with multiple lines for hemodynamic monitoring (e.g., arterial line, central venous catheter, pulmonary artery catheter) and medication administration. High-dose immunosuppressive therapy is initiated immediately to prevent acute rejection. The focus is on stabilizing hemodynamics, managing pain, and monitoring for primary graft dysfunction. Chest tubes will be in place for drainage.
  
  \item **First 1-2 Weeks (Hospital Ward):** Once stable, usually within a few days, the patient is transferred to a specialized transplant unit. Mechanical ventilation is weaned, and lines are removed. Early mobilization, often with the assistance of physical therapists, begins. The focus shifts to continued recovery, medication education (including diet restrictions like avoiding grapefruit), infection prevention, and wound care. Diet progresses from clear liquids to a cardiac-friendly diet. The length of hospital stay varies but typically ranges from two to four weeks.
  
  \item **First 1-3 Months (Outpatient):** After discharge, patients face a lifelong commitment to immunosuppression, frequent clinic visits (weekly, then bi-weekly), and a vigilant approach to monitoring for rejection and infection. Regular endomyocardial biopsies and echocardiograms are performed. Cardiac rehabilitation is crucial for regaining strength and functional capacity. Activity restrictions are gradually lifted, but heavy lifting and strenuous activities are avoided initially.
  
  \item **Long-term Recovery (Beyond 3 Months):** Patients transition to less frequent clinic visits (monthly, then quarterly, then annually). Lifelong immunosuppression continues, with careful monitoring of drug levels and side effects. Annual surveillance for coronary allograft vasculopathy (CAV) with coronary angiography is common. Most patients can return to work, exercise, and normal daily activities, albeit with continued adherence to their medical regimen and lifestyle modifications.
\end{itemize}

\section*{Surgical Findings \& Pathology}
During the heart transplant procedure, the surgeon documents the condition of the recipient's native heart, noting any significant hypertrophy, dilation, fibrosis, valvular abnormalities, or presence of thrombus. The extent of pericardial adhesions, particularly from previous cardiac surgeries, is also noted as it can complicate the cardiectomy. Signs of pulmonary hypertension in the recipient's pulmonary artery may also be observed.

The resected native heart is sent for pathological examination. The pathology report confirms the underlying etiology of end-stage heart failure, such as ischemic cardiomyopathy (evidenced by myocardial scarring), dilated cardiomyopathy (marked by ventricular dilation and thinning), hypertrophic cardiomyopathy (showing myocardial thickening and disarray), or restrictive cardiomyopathy. It details the extent of myocardial fibrosis, necrosis, inflammation, and any valvular pathology or intracardiac thrombi. Post-transplant, surveillance endomyocardial biopsies of the donor heart are the gold standard for diagnosing acute cellular rejection, graded according to the International Society for Heart and Lung Transplantation (ISHLT) scale (e.g., 0R for no rejection, 1R for mild, 2R for moderate, 3R for severe). Antibody-mediated rejection is assessed through immunofluorescence and immunohistochemistry (e.g., C4d staining). These biopsies are crucial for guiding immunosuppressive therapy adjustments.

\section*{Expected Postoperative Course}
A normal, successful postoperative course following heart transplantation is characterized by a significant and sustained improvement in the patient's overall health and functional status. Patients typically experience a dramatic resolution of their pre-transplant heart failure symptoms, such as severe dyspnea, fatigue, and peripheral edema, often within weeks to months. The transplanted heart demonstrates normal cardiac output and stable hemodynamics, confirmed by echocardiography showing a normal left ventricular ejection fraction and absence of significant valvular dysfunction. Pain from the sternotomy incision gradually resolves over several weeks. Patients are expected to regain strength and endurance through cardiac rehabilitation, progressively increasing their activity levels and often returning to a functional status of NYHA Class I or II. The absence of significant acute rejection episodes or major opportunistic infections, coupled with stable immunosuppressant drug levels, signifies a well-managed course. Ultimately, the goal is for the patient to resume many normal daily activities, including work, exercise, and social engagement, with a renewed quality of life, albeit with a lifelong commitment to medical follow-up and adherence to their treatment regimen.

\section*{Postoperative Care \& Follow-up}
Lifelong, rigorous follow-up is essential after heart transplantation to monitor for rejection, infection, and other complications.

\begin{itemize}
  \item **Regular Clinic Visits:** Frequent visits initially (weekly, then monthly) transitioning to less frequent (quarterly, then annually) as the patient stabilizes.
  
  \item **Endomyocardial Biopsies:** Scheduled biopsies are performed frequently in the first year (e.g., weekly, then bi-weekly, then monthly) and then less often, or as indicated for suspected rejection.
  
  \item **Echocardiograms:** Regular echocardiograms to assess cardiac function and detect early signs of dysfunction.
  
  \item **Coronary Angiography:** Typically performed annually starting one year post-transplant to screen for coronary allograft vasculopathy (CAV).
  
  \item **Blood Tests:** Frequent monitoring of immunosuppressant drug levels, renal and hepatic function, electrolytes, complete blood count, and viral titers (e.g., CMV, EBV).
  
  \item **Physical Therapy and Cardiac Rehabilitation:** Essential for regaining strength, endurance, and overall functional capacity.
  
  \item **Skin Cancer Screening:** Regular dermatological examinations due to increased risk of skin malignancies from immunosuppression.
  
  \item **Vaccinations:** Annual influenza vaccine and other recommended vaccinations (e.g., pneumococcal, tetanus) are crucial, often with specific guidelines for immunosuppressed patients.
  
  \item **Education on Warning Signs:** Patients are educated to immediately report symptoms such as fever, chills, shortness of breath, sudden weight gain, swelling, chest pain, or new arrhythmias, which could indicate rejection or infection.
\end{itemize}

Adherence to this comprehensive follow-up schedule is paramount for long-term success and survival.

\section*{Epidemiology}
Heart transplantation remains a relatively rare but life-saving procedure, constrained primarily by the limited availability of donor organs. Globally, approximately 3,500 to 4,000 heart transplants are performed annually, with the United States accounting for a significant proportion of these (around 3,800 in 2022). The incidence of end-stage heart failure, the primary indication for transplant, is much higher, highlighting the disparity between need and availability. Heart transplant recipients are predominantly adults, with the peak age range typically between 50 and 70 years, though pediatric heart transplantation is also performed. Males constitute a larger proportion of recipients compared to females, often reflecting the higher incidence of ischemic cardiomyopathy in men. The most common indications for transplantation are ischemic cardiomyopathy and dilated cardiomyopathy, together accounting for over 70\% of cases. Despite the inherent risks, 1-year survival rates are excellent, typically ranging from 85-90\%, with 5-year survival rates around 70-75\%, and 10-year survival rates approaching 50-55\%. Morbidity remains high, largely due to complications arising from lifelong immunosuppression, including opportunistic infections, renal dysfunction, hypertension, diabetes, and the development of coronary allograft vasculopathy.

\section*{Outcomes \& Prognosis}
The outcomes of heart transplantation have dramatically improved since its inception, offering a significantly prolonged and enhanced quality of life for patients with end-stage heart failure. Typical success rates for survival are high, with most centers reporting 1-year survival rates exceeding 85\% and median survival now approaching 12-13 years. Functional outcomes are generally excellent, with the vast majority of recipients achieving a significant improvement in their functional capacity, often returning to NYHA Class I or II, enabling them to resume work, exercise, and normal daily activities. While the original heart disease does not recur, the primary challenges to long-term prognosis are chronic rejection, particularly coronary allograft vasculopathy (CAV), and complications related to lifelong immunosuppression, such as infections, renal dysfunction, and malignancy. Prognostic factors influencing long-term survival include donor age and quality, recipient age and comorbidities (e.g., pre-existing diabetes, renal impairment), the presence of pre-transplant pulmonary hypertension, early graft function, and, critically, the patient's adherence to their immunosuppressive regimen and follow-up care. Despite these challenges, heart transplantation remains the gold standard for end-stage heart failure, offering a robust prognosis for carefully selected patients.

\section*{References}
\begin{enumerate}
  \item MedlinePlus. Heart transplant. U.S. National Library of Medicine; 2024. Available from: https://medlineplus.gov/hearttransplant.html
  
  \item Braunwald E, Zipes DP, Libby P, Bonow RO, Mann DL, Tomaselli GF. Braunwald's Heart Disease: A Textbook of Cardiovascular Medicine. 12th ed. Elsevier; 2022.
  
  \item Sabiston DC, Townsend CM. Sabiston Textbook of Surgery: The Biological Basis of Modern Surgical Practice. 21st ed. Elsevier; 2021.
  
  \item International Society for Heart and Lung Transplantation (ISHLT). ISHLT Guidelines for the Care of Heart Transplant Recipients. J Heart Lung Transplant. 2024.
\end{enumerate}

\clearpage